\chapter{Introduction}\label{ch:introduction}

% From Paper 1: "Entity Multiplexing Through Activation Strength: An Application of Sparse Autoencoders to Reinforcement Learning"
Understanding how deep learning models function is currently an unsolved problem in the field of machine learning. This state of affairs ensures that highly reliable, controllable and understandable models remain out of reach. This is particularly the case in the field of reinforcement learning (RL), which has been somewhat neglected by the techniques of mechanistic interpretability (MI) whose focus has largely been on Large Language Models (LLMs). With techniques from RL being applied to frontier models to a greater extent, it is our belief that applying MI techniques to smaller AI systems trained using RL can yield insights that might lead to a more general understanding of fundamental aspects of neural network operations.

% From Paper 2: "Discovering Goals in A Deep Reinforcement Learning Model by Studying Steering Vectors"
As RL techniques become more effective and are deployed more widely, it will become increasingly necessary to provide safety guarantees and explanations behind critical trade-offs that RL policies are making so that we can deploy these systems with confidence that they will work as we expect\citep{shakya_2023_reinforcement}.

Existing techniques are being used to better understand RL models\citep{atrey_2020_exploratory}, but many techniques developed for mechanistic interpretability elsewhere have not yet seen much application in reinforcement learning (\cite{hilton_2020_understanding} being a notable exception). There is an opportunity to explore novel techniques for interpreting other categories of AI systems with RL, as RL environments usually contain relatively simple goals and easily determined failure modes that allow us to test whether changes we make evoke a behavior we like. For example, the approach of steering Large Language Models (LLMs) using activation steering in\cite{turner_2024_activation} was inspired by the work done to understand and control a maze solving agent in the Procgen Maze environment\cite{mini_2023_understanding}.

Our investigation seeks to build on the work of\cite{mini_2023_understanding}, who found key parts of the model that could be used to control and steer behavior. Our contribution lies in extending their work to an environment that involves multiple competing entities that an agent needs to reach in sequence. To this end, we chose the Procgen Heist environment as the target of our research. By exploring how models trained in the Procgen Heist environment\citep{cobbe_2020_leveraging} approach multiobjective maze solving, we aim to understand how neural networks represent and pursue multiple sequential goals.

Initially, extracting meaningful insights from this model proved challenging due to the high degree of polysemanticity within the network that made it challenging to draw strong conclusions about how it worked. To address this, we train Sparse AutoEncoders (SAEs) on the CNN layers of our network in order to disentangle the functionality of the model. SAEs are a type of unsupervised learning technique that can be used to break down the most critical functions of a model into a more interpretable set of sparse latents by training an overcomplete basis on the original activations of a given layer of the network.

% Outline of the thesis structure (placeholder)
This thesis is structured as follows: Chapter\ref{ch:literature} provides a review of the relevant literature in reinforcement learning and machine learning interpretability. Chapter\ref{ch:methodology} details the experimental setup, including the environment, model architecture, and the interpretability techniques employed. Chapter\ref{ch:results} presents the findings from our experiments, including feature visualizations, steering results, and ablation studies. Chapter\ref{ch:discussion} discusses the implications of these findings, and Chapter\ref{ch:conclusion} concludes the thesis and suggests avenues for future research.

\section{Background and Motivation}\label{sec:background}

Reinforcement learning (RL) is a field of machine learning that has the unique promise of delivering superhuman performance in a wide range of environments\cite{Silver2016}. However, the black box nature of these models has made it difficult to understand how they achieve these results. In recent years, RL has seen a surge in interest from the ML community, due to its application in LLMs and other frontier models\cite{ouyang2022training, bai2022training}. This has led to significant advances in performance in these models, but also presents the potential for these highly capable systems to manifest unexpected behavior that is unique to the RL setting.  

Machine learning as a field is composed primarily of three different approaches to training models. These are supervised learning, unsupervised learning, and reinforcement learning. Supervised learning requires large amounts of carefully labeled data to allow models to learn the underlying patterns that connect the input to the specific label. Unsupervised learning has no labeled data and simply has the model learn the structures in the data through identifying patterns in the data itself.  

Reinforcement learning differs in that it focuses on the idea of an environment that is the source of truth for the agent. By providing actions to the environment based on its policy, the agent receives new information in the form of observations from the environment, as well as reward information that provides a signal about whether it is achieving the goals that the designers of the model and environment intended. This is typically done by using the reward as a key variable in the loss function, with higher rewards being linked to lower loss.

This means that with relatively low amounts of human effort in shaping environment behaviour, agents trained with RL can achieve performance that can far surpass human capabilities in a domain, as was demonstrated by systems such as\citep{Silver2016} AlphaGo. In that case the agent was trained through self play whereby it played against a similarly skilled version of itself, and would learn to improve on its mistakes based on whether the moves it made in a game led to victory or defeat.

Modern RL often employs deep neural networks in order to achieve some of its most remarkable results, but this leads to the issue of the lack of understandability of how those results are actually achieved due to the black box nature of modern neural networks.

The field of mechanistic interpretability (MI) exists in order to help us peer into and decipher these black boxes so that we can make use of these models with greater confidence and reliability, and to allow us to be proactive about finding issues in models' behaviour before critical issues arise. 

MI as a field aims to answer the question of how neural networks actually function at a very granular level. In the ideal case it would make it possible to reverse engineer a given neural network into code through a complete knowledge of how its parts work. Having this understanding would give us a broad box of tools for helping us to make the neural networks safer, and more reliable.

MI broadly refers to a number of approaches that we can use to gain insight into neural network operations, from visualising neurons and learned filters in CNNs, to the training of specialised neural networks that help us better understand the operations of a particular part of a network. 

The ultimate goal of mechanistic interpretability has become to find ways to ensure the safety of AI systems, and this has become particularly critical with frontier agentic AI that has the ability to perform long horizon tasks that could involve taking actions that are highly impactful to humans. It is critical to be able to understand and prevent unwanted intent from developing in such systems, and to do this it is necessary to understand what goals the system may have. 

For a long time, many in the field questioned whether models had an internal model such that it would even make it coherent to discuss whether a neural network could have things such as 'intent' or 'goals'. However recent work has provided compelling evidence that neural networks can indeed learn meaningful internal models and conceptual abstractions from data. For instance, studies on language models trained only on the moves made in a game such as Othello, with no explanation of the rules, revealed that these models spontaneously developed internal representations of the game state without explicit supervision of the rules \citep{li2022emergent}. 

Other research has shown that such models naturally learn geographic mappings of the world related to locations that can be retrieved directly from the activations of the network, effectively building simplified 'world models' \citep{gurnee2023language}. Similarly, analyses of agents trained on complex games like chess have demonstrated that these systems acquire human-understandable strategic concepts, such as piece values or tactical motifs, purely from self-play \citep{mcgrathAcquisitionChessKnowledge2022}. These findings suggest that the internal workings of these models are richer than previously thought and lend support to the idea that they may plausibly learn to represent goals that can be directly extracted from the model and understood.

It therefore seems plausible that something as fundamental as goals would be likely to share some properties in the way it is encoded across models. With that hope in mind, this work seeks to reveal some of these properties to inform investigations across future models.

However not all models which seek to target goals are good candidates for this research, as it may easily be the case that the task is simple enough to solve without the need to develop a sufficiently complex mechanism to represent abstract concepts like goals. The procgen environments stand out as good environments for this reason, as they procedurally generate environments for the agent to solve. This means that the agent cannot simply overfit to the data but has to learn general patterns in order to score well on the environment, which would necessitate developing a complex representation of goals within the model.

The choice to use Procgen Heist in particular stemmed from the desire to explore how multiple goals might be represented in the model, as significant work has already shown the presence of "cheese representations" with the Procgen Maze environment\citep{mini_2023_understanding}. Because advanced agentic AI would be operating environments with potentially multiple goals, the question of whether there was a substantial difference in the representation of these goals was an essential one.

However, in making the environment more complex the model too became more complex, and the problem of polysemantic neurons became a significant hurdle to understanding the model's operations. Specifically whether the model's representations of goals in the base model represented a "true" representation or whether it was simply the results of polysemanticity in the network. 

Polysemanticity refers to the case of a single neuron having multiple coherent feature representations encoded simultaneously, due to it needing to encode many such features but not having sufficient neurons to represent them all. In such cases the model will compress multiple features into a single neuron, typically encoding features that are unlikely to activate strongly simultaneously in practice to avoid causing harmful interference. 

This led to our use of Sparse AutoEncoders (SAEs) to confirm whether neurons in the network would naturally tend to represent the different entities using different neurons, or whether there were single 'current goal' channels in the network that would naturally come to exist. 

An SAE is a model trained with an unsupervised approach to simply optimally reconstruct the outputs of a given layer in the network. Typically it is trained with a loss function that minimises the reconstruction loss, and the L1 norm of the network, which means that it produces the smallest possible activations. Typically the SAE has a hidden layer between an encoder and decoder which are the primary parameters of the network, and the size of the hidden layer is normally based on the size of the original layer multiplied by some expansion factor. This means that the SAE learns an overcomplete basis of the original layer, which allows for the separation of polysemantic features into separate neurons in the network, which can then be analysed in isolation.

Being able to understand how these models work is a key step towards building more reliable, controllable and understandable models. Specifically, we are interested in understanding how these models represent and pursue goals in complex, multi-objective environments that are a step beyond simple toy environments. Goals, specifically are something that advanced agentic AI would share with even simple models, and so greater knowledge on the underlying principles behind this universal concept may provide us with fundamental insights that generalise beyond toy environments.



\section{Problem Statement}\label{sec:problem}


\section{Research Objectives}
\label{sec:objectives}
% List your research aims and objectives
% These should be specific, measurable, and achievable

\section{Research Questions}
\label{sec:questions}
% State the specific questions your research seeks to answer

\section{Scope and Limitations}
\label{sec:scope}
% Define the boundaries of your research
% Acknowledge what is beyond the scope of your study

\section{Significance of the Study}
\label{sec:significance}
% Explain the theoretical and practical contributions of your research
% Discuss the potential impact of your findings

\section{Thesis Structure}
\label{sec:structure}
% Provide an overview of the chapters that follow
% Briefly describe the content of each chapter